%% ============================================================
%%  Introducción
%%  Duenio inicial: claude-1
%%  Reglas: ver ../../CLAUDE.md
%%  No editar si TASKS.md marca este archivo IN_PROGRESS por otro agente.
%% ============================================================

\section{Introducción}

\subsection{Contexto}

En numerosos problemas reales, los modelos profundos --en particular los basados en arquitecturas Transformer-- alcanzan un desempeño predictivo superior al de los modelos matemáticos clásicos. Sin embargo, su naturaleza de caja negra impide entender por qué producen una salida dada, lo cual limita su adopción en dominios donde la interpretabilidad es crítica. La limitación es especialmente aguda en las ciencias de la Tierra, donde el objetivo no se agota en predecir: se espera que el modelo genere conocimiento sobre el proceso subyacente \cite{reichstein2019deep}. Un modelo que predice bien la dinámica de un cultivo a partir de índices espectrales, pero que no permite formular hipótesis sobre las relaciones entre esos índices, tiene un valor científico acotado.

La mayoría de los métodos de explicabilidad (XAI) disponibles se limita a tres operaciones: atribuir importancia marginal a variables de entrada mediante valores de Shapley o gradientes integrados \cite{medtext2026attribution}; visualizar la relevancia sobre la atención como un mapa continuo, con o sin corrección por propagación de relevancia y gradientes \cite{chefer2021transformer,abnar2020quantifying}; o inyectar en la arquitectura una estructura relacional conocida de antemano, de modo que el grafo es insumo y no resultado \cite{zhao2026causalguided}. Ninguna de las tres convierte la representación latente en una hipótesis relacional compacta que pueda compararse, validarse y generalizarse entre dominios: las dos primeras explican el modelo o la entrada, y la tercera parte de una estructura dada, que es precisamente aquello que este trabajo busca descubrir.

Existe, no obstante, un antecedente directo. La conversión de la representación interna de un Transformer en estructura de grafo ya ha sido ensayada: se han construido grafos causales del flujo de información entre cabezas de atención de un Transformer de series temporales ya entrenado, interviniendo sus activaciones mediante parcheo \cite{mechinterp2025ts}, y existen arquitecturas que aprenden la topología entre variables como componente interno del modelo \cite{cai2024msgnet,pipeline2026causalgraph}. La diferencia que este trabajo introduce no está en la extracción sino en el régimen de validación y en el mecanismo de acceso al modelo. Un grafo aprendido dentro del modelo existe porque fue entrenado para existir, y auditarlo exige reentrenar, lo que lo vuelve inaplicable a un modelo ya desplegado; un grafo obtenido por intervención exige parchear activaciones y acceso a los estados internos del modelo; y un grafo extraído de una única corrida de entrenamiento no informa sobre su propia reproducibilidad. El framework propuesto se ubica una capa después: toma como entrada la matriz que esos modelos ya producen, sin intervenir su cómputo interno, y traslada la carga de la prueba desde el desempeño predictivo hacia la validación estadística e intervencional --sobre la hipótesis ya extraída, no sobre el modelo-- de cada relación. Frente a los enfoques que inyectan un grafo causal preexistente, no impone estructura relacional previa; frente a los que la aprenden de extremo a extremo o la obtienen por intervención, no exige reentrenar ni acceso a activaciones internas para auditar.

\subsection{El problema}

El problema central es que la atención de un Transformer no es un objeto estable. Pequeñas variaciones en la semilla de inicialización producen matrices de atención numéricamente distintas aun cuando el desempeño predictivo sea equivalente, de modo que las matrices obtenidas en entrenamientos independientes no son intercambiables entre sí. Esto invalida cualquier explicación derivada de un único entrenamiento y obliga a un enfoque de consenso sobre un ensamble de modelos, en el que la unidad de análisis deja de ser la matriz de una corrida y pasa a ser la frecuencia con que una relación reaparece.

A esa inestabilidad se suman dos problemas documentados. Primero, la fragilidad de las explicaciones post-hoc ha sido demostrada de forma adversarial: es posible construir modelos que ocultan sistemáticamente su comportamiento real a los métodos de atribución basados en perturbación \cite{slack2020fooling}. Segundo, un modelo con métricas de desempeño excelentes puede estar apoyándose en atajos espurios y no en el fenómeno de interés, efecto que solo se hace visible al inspeccionar la estructura de la evidencia y no el error agregado \cite{lapuschkin2019unmasking}.

Un tercer obstáculo es metodológico: la evaluación de explicaciones carece de consenso métrico. Una evaluación sistemática sobre Vision Transformers concluye que las métricas que dicen medir una misma propiedad --fidelidad, coherencia, robustez, complejidad-- convergen y correlacionan mínimamente entre sí, de modo que el orden de mérito de un método depende de cuál métrica se elija \cite{vitexplain2024evaluation}. La comparación directa de dos explicadores igualmente fieles --SHAP y gradientes integrados, ambos en torno a 0,22 de fidelidad-- muestra además que solo correlacionan de forma moderada entre sí: coincidir en la fidelidad no implica coincidir en qué se señala como relevante \cite{medtext2026attribution}. Este trabajo responde a esa inconsistencia con una variante: no se busca un número que certifique la explicación, ni un escalar que agregue los criterios, sino un conjunto de criterios que deban converger y cuya discordancia sea, ella misma, informativa.

El problema puede entonces enunciarse en una frase: no existe un procedimiento que permita decidir, con criterios declarados de antemano, cuándo una relación entre variables extraída de la representación latente de un modelo profundo constituye una hipótesis sobre el fenómeno y cuándo es un artefacto del entrenamiento.

\subsection{Justificación: por qué este estudio y qué aporta}

La pregunta que motiva el trabajo no es si es posible extraer un grafo desde una matriz de atención --lo es, y con procedimientos simples-- sino si ese grafo puede sostenerse como afirmación científica. Esa distinción es la que hoy falta en la literatura, y su ausencia tiene un costo concreto: las estructuras relacionales que se publican como hallazgo se reportan sobre una única corrida de entrenamiento, sin frecuencia de persistencia entre semillas, sin margen respecto del corte de discretización y sin contraste contra un modelo nulo. En consecuencia, no es posible distinguir una relación aprendida del fenómeno de una coincidencia de inicialización.

El aporte que este trabajo reclama es, por tanto, doble y de naturaleza metodológica. Por un lado, provee el objeto: una hipótesis relacional discreta, dirigida, definida sobre variables con nombre y de tamaño acotado, que puede enunciarse fuera del modelo que la produjo. Por otro, y sobre todo, provee el protocolo que permite rechazarla. Un mapa de relevancia continuo no admite esa operación, porque no ofrece una unidad discreta que eliminar; un grafo impuesto a priori tampoco, porque no existe hipótesis nula respecto de la cual pueda resultar improbable. Esta combinación --entregar simultáneamente una estructura y el criterio explícito para descartarla-- es la que la literatura revisada no reúne.

Tres razones adicionales justifican emprenderlo ahora. La primera es de oportunidad metodológica: la evidencia de que las métricas individuales de XAI son mutuamente inconsistentes \cite{vitexplain2024evaluation,medtext2026attribution} y de que las basadas en perturbación son manipulables \cite{slack2020fooling} está establecida, pero no se ha traducido todavía en un protocolo de validación de estructura relacional. La segunda es de oportunidad arquitectónica: tratar cada variable como token, en lugar de cada posición espacial o temporal, se ha consolidado como decisión efectiva y no meramente conveniente \cite{liu2024itransformer,cai2024msgnet}, lo que hace que hoy existan matrices de atención entre variables disponibles en modelos publicados de múltiples dominios. La tercera es de alcance: al depender únicamente de las propiedades algebraicas de la matriz de entrada y no de su procedencia, el framework es reutilizable por terceros sobre modelos que no fueron diseñados para él.

\subsection{La propuesta}

Se propone un framework estructurado en cuatro etapas:

\begin{enumerate}
\item Tomar la representación latente de un modelo profundo arbitrario, ya entrenado.
\item Transformarla en un grafo dirigido ponderado ($\Phi: A \to G$).
\item Extraer una hipótesis relacional compacta ($\Psi: G \to H$).
\item Someter esa hipótesis a validación mediante estabilidad estadística y perturbación dirigida ($F: H \times M \to \mathbb{R}$).
\end{enumerate}

El cuarto paso no se resuelve con una métrica. Se plantea como una cadena de evidencia acumulativa de cinco eslabones --estabilidad continua de $A$; estabilidad estructural bajo distintos esquemas de discretización; margen Top-K contra un modelo nulo con corrección por comparaciones múltiples; fidelidad predictiva por ablación de aristas; y especificidad, robustez y compacidad--, complementada con una triangulación de cinco paradigmas ortogonales: probabilístico, intervencional/causal, geométrico-topológico, teoría de la información y baselines deterministas. Bajo este criterio, una arista solo se considera parte del núcleo relacional si sobrevive a eslabones sucesivos y es respaldada por paradigmas que no comparten supuestos.

Cabe precisar el estado de avance: los dos primeros eslabones cuentan con evidencia preliminar sobre un prototipo del caso de estudio. Entrenados 50 modelos con semillas distintas, la arista dominante persiste en el 76--78\% de ellos, con un Jaccard promedio de 0.195--0.204 entre las matrices de atención resultantes --un núcleo estable con periferia variable--, y un análisis de sensibilidad al variar el umbral de discretización ($K\in\{3,5,7,10\}$) no encuentra un punto crítico, lo que sugiere que también la periferia se degrada de forma gradual y no abrupta. Los eslabones 3 a 5 --consenso estadístico frente a un modelo nulo, fidelidad predictiva por ablación y especificidad/robustez-- y el resto de la triangulación de paradigmas permanecen, en cambio, como diseño metodológico propuesto, aún no ejecutado. La interpretabilidad, en este marco, no busca una única verdad estática, sino una hipótesis relacional de consenso robusta frente a la variabilidad estocástica del entrenamiento.
